\documentclass[10pt]{article}
\usepackage[margin=1in]{geometry}
\usepackage{booktabs,hyperref}
\title{Supplement to Auditable Empirical Evaluation of Machine Unlearning}
\author{Matthew Looney (Magnexis)}
\date{Research Release R001 --- P001}
\begin{document}
\maketitle
\section{Scope}
This supplement inventories the reproducibility contract for the retrospective P001 pilot.
It is not evidence for C001 or for general unlearning effectiveness.
\section{Model and optimization}
The complete architecture, initialization, optimizer, learning rate, epoch count, batch
size, and method configurations are preserved in the P001 configuration and raw record.
The release verifier checks their hashes rather than restating them manually here.
\section{Seeds and manifests}
The exact-retraining seeds are 2001--2015. The release includes machine-readable deletion
and retained indices and verifies that they are disjoint and jointly reconstruct the frozen
training split.
\section{Dataset integrity}
The reduced pilot contains 9,000 training, 1,000 validation, and 2,500 test examples. The
canonical post-scaling dataset SHA-256 is
\texttt{1ad69a47dbce37a98443e5e5628a031b5910be879d3589e945a1519f805aa28c}.
\section{Quality gates and storage}
Schema, hashes, split integrity, deleted-set exclusion, finite metrics, checkpoint and
prediction references, raw immutability, and figure provenance are checked automatically.
Raw evidence and derived summaries remain separate.
\section{Exact/exact and convergence calculations}
Exact/exact distances characterize stochastic retraining variability. The frozen pilot
diagnostic compares estimates based on the first ten and all fifteen exact seeds. Pairwise
distances are not treated as independent experimental units in later power simulations.
\section{Approximate methods}
P001 executed fine-tuning, naïve gradient ascent, and SCRUB. The first two are simple
engineering baselines. Because the deletion condition was retrospectively assay-insensitive,
their output proximity is not interpreted as successful forgetting.
\section{Runtime and environment}
Runtime metadata records WSL2, Python 3.12.3, PyTorch 2.9.1 with ROCm 7.2, and the AMD RX
7800 XT. Runtime numbers are configuration-specific descriptive measurements.
\section{Failure and amendment records}
The failure registry, corrections file, negative-result record NR001, decision log, and
registration drafts are included or referenced by the release inventory.
\section{Reproduction protocol}
Run \texttt{python scripts/verify\_r001\_release.py} for Level A artifact verification.
Level B retraining uses the frozen P001 configuration and must create a new output record;
it must not overwrite the archived experiment.
\section{Artifact inventory and timeline}
The release checksum manifest is authoritative. HISTORY.md and research/PROJECT\_TIMELINE.md
distinguish retrospective P001 work from the not-yet-executed prospective C001 study.
\end{document}
